\input{preamble}

\begin{document}

\title{Supplement to ``Exact Marginal--Regret Frontiers\\for Signed-Prefix Compression'':\\
Proofs and Formal Declarations}

\author{\IEEEauthorblockN{Vu Phan An Nguyen}
\IEEEauthorblockA{\textit{Rider University} \\
Lawrenceville, NJ, USA}}

\maketitle

\begin{abstract}
This supplement gives, for each result of the main paper, the proof architecture in
readable mathematics and every Lean~4 declaration it rests on, grouped by module. It also
records the status of declarations that the main text does not use: superseded routes,
arithmetic lemmas whose inputs are not instantiated, and one withdrawn theorem whose domain
is empty. Section numbers S1--S7 follow the main text; theorems are named as in the main
text.
\end{abstract}

\section*{Conventions}

All declarations live in \texttt{work/p13\_conjectures}. A module is accepted when it
compiles with exit status $0$, contains no \texttt{sorry} and no \texttt{native\_decide}
outside comments, and the axioms of every audited theorem lie within
\{\texttt{propext}, \texttt{Classical.choice}, \texttt{Quot.sound}\}. Acceptance is checked by
a frozen gate that copies every module and its local imports into a fresh directory and
builds from an empty cache (S7). The paper's objects are \lean{gFamily} ($g_{m}$),
\lean{regret} ($r_{k}$), \lean{FeasibleAt} (the marginal band of the main text, read on
the dropped prefix and equivalently on the kept complement, quantified over budgets
$k\le m$), \lean{value} ($V_{m}(\eta)$), \lean{freeValue}
($V^{\mathrm{free}}_{m}$) and \lean{trueValue} (the physical optimum $V(g';\eta)$).

\section{S1. The family and its baseline}

\paragraph{The predicates do not pin the family} Much of the development is stated for the
predicates ``weights in $\{-3,-1,1,3\}$, total $0$'' (even) or ``total $-1$'' (odd). Those
do not force two heavy items: $(1,1,-1,-1)$ satisfies the even predicate with none. The
endpoint arguments use $|\mathrm{heavySet}|=2$, so $g_{m}$ is written down, its total is
computed by peeling the two heavy endpoints and splitting the middle at the sign change, and
its heavy set is shown to be exactly the two endpoints.

\leanbox{\lean{gFamily}, \lean{gFamily\_total}, \lean{gFamily\_isEvenFamily},
\lean{gFamily\_isOddFamily}, \lean{gFamily\_heavySet}, \lean{gFamily\_heavy\_card},
\lean{isEvenFamily\_not\_forcing\_two\_heavy} in \lean{LeanPaperCFamilyV1};
\lean{card\_weightClass} in \lean{LeanPaperCEndpointLawV1}; \lean{sum\_fn\_gFamily} in
\lean{LeanPaperCRelaxedJoinV1}; \lean{normMass\_gFamily} in \lean{LeanPaperCRatioJoinV1}}

\paragraph{The baseline} Every weight is odd, so $S_{k}$ has the parity of $k$ and
$|S_{k}|\ge k\bmod2$ for every ranking; using $k\bmod2$ therefore never understates regret.
The other side: for each interior budget an explicit subset attains $k\bmod2$ (light items
only when $k\le m-2$; at $k=m-1$ every item but one light negative). Hence
$b_{k}(g_{m})=k\bmod2$ exactly, and the physical regret with the baseline recomputed from
its definition equals the parity excess on $g_{m}$. This is a statement about subsets, one
budget at a time; a single ranking meeting $k\bmod2$ at every budget is neither claimed nor
needed.

\leanbox{\lean{prefixSum\_parity} in \lean{LeanPaperCRankingV1};
\lean{baseline\_le\_abs\_prefix} in \lean{LeanPaperCValueV1}; \lean{attainment},
\lean{familyBaselineIdentification\_holds},
\lean{family\_regret\_identification\_unconditional} in \lean{LeanPaperCBaselineAttainV1};
\lean{trueValue\_family} in \lean{LeanPaperCStabilityFamilyV1}}

\paragraph{The prefix is the dropped set} In the physical model the loss of keeping a set $K$
is half the oscillation of the weights outside it,
$\mathrm{selectedRadius}(g,K)=\tfrac12|\sum_{i\notin K}g_{i}|$. The regret of the main
text applies this with $K$ the complement of the prefix, so the first $k$ items are the ones
dropped, and the kept set has size $m-k$. The best physical loss at that size is exactly
half the baseline, and dividing by $\|g\|_{1}/2$ gives the normalised regret. On
$g=(-3,-1,3)$ at $k=1$ the two readings differ ($\tfrac27$ against $\tfrac47$), so the
direction is not a convention.

\leanbox{\lean{physical\_radius}, \lean{physical\_minimum}, \lean{physical\_normalized\_regret}
in \lean{LeanPaperCStabilityEndgameV1}}

\paragraph{The feasible set is inhabited} An earlier version of \lean{FeasibleAt}
quantified the band over all $k\in\mathbb{N}$. For $k>m+\eta$ every item is in the
prefix and $|1-k/m|>\eta/m$, so no law was feasible at any finite $\eta$ and every value
statement was vacuous. The quantifier
is bounded by $k\le m$, and the set is proved non-empty for every $\eta\ge0$ on the
permutation laws the value ranges over.

\leanbox{\lean{FeasibleAt}, \lean{feasible\_mono} in \lean{LeanPaperCValueSetV1};
\lean{allPerm\_feasible\_nonempty} in \lean{LeanPaperCLawPushV1}}

\section{S2. Zero-band prices (Theorem ``Zero band'')}

\paragraph{Even $m$, upper side} The extra ranking $(-3,+1,-1,\dots,+1,-1,+3)$ has prefix
sum $-3$ at odd budgets and $-2$ at even ones, so its excess is exactly $2$ at every
interior budget, and for even $m$ its total is zero, so it is a ranking of the family. The
$m-1$ insertion rankings carry a total excess of $2$ at every budget (each a single spike,
the spikes covering the budgets once). Mixing all $m$ uniformly gives $4/m$ at every
budget, i.e.\ $\frac{4}{mD}$. A sign flip at ratio $1/2$ (free, since excess is built from
$|S_{k}|$) and uniform relabelling within each weight class make every per-item marginal
exactly $k/m$.

\leanbox{\lean{flatEven\_prefix}, \lean{flatEven\_excess}, \lean{flatEven\_total},
\lean{flatEven\_weights} in \lean{LeanPaperCEvenV1}; \lean{insertion\_total\_excess},
\lean{evenMixture\_value}, \lean{evenMixture\_value\_at\_D},
\lean{evenMixture\_value\_constant} in \lean{LeanPaperCMixV1}; \lean{excess\_neg} in
\lean{LeanPaperCRankingV1}; \lean{uniform\_band\_zero} in \lean{LeanPaperCUniformV1}}

\paragraph{Even $m$, lower side} For even $m$ the prefix of length $m-1$ is minus the last
item, so $r_{1}+r_{m-1}=(2/D)(\mathrm{heavy}(\pi_{1})+\mathrm{heavy}(\pi_{m}))$. At $\eta=0$
each end is heavy with probability $2/m$, so the objective is at least $\frac{4}{mD}$ for every
feasible law.

\leanbox{\lean{two\_budget\_identity} in \lean{LeanPaperCRankingV1};
\lean{value\_ge\_at\_uniform\_at\_D} in \lean{LeanPaperCLowerZeroV1};
\lean{even\_value\_zero} in \lean{LeanPaperCEndgameV1}}

\paragraph{Odd $m$: the window sums} The upper side is the uniform law on the cyclic shifts
of $(+3,-3,-1,+1,\dots,-1)$. Working on $\mathbb{N}$ with a range split rather than in
$\mathbb{Z}_{m}$, the window sum starting at $a$ of length $k$ is
$W(a,k)=(P(\min(a+k,m))-P(a))+P(a+k-m)$, where $P$ is the ordinary prefix sequence, with
$P_{0}=0$, $P_{1}=3$ and, for $j\ge2$, $0$ at even $j$ and $-1$ at odd $j$. For odd $m\ge3$
and $0<k<m$,
\[
\sum_{a<m}|W(a,k)|=\begin{cases}m+4,&k\text{ odd},\\ k+4,&k\text{ even},\end{cases}
\]
so the value at budget $k$ is $4/m$ at odd $k$ and $(k+4)/m$ at even $k$. The maximum is
$(m+3)/m$, at $k=m-1$ only. Odd budgets carry the \emph{larger} raw sum and the
\emph{smaller} value, because $b_{k}=1$ removes $m/m$ there. The base sequence is a
rearrangement of $g_{m}$: its weight counts agree with those of $g_{m}$ at every real
weight.

\leanbox{\lean{win\_eq}, \lean{preN\_pattern}, \lean{preN\_alt\_sum},
\lean{win\_abs\_sum\_odd}, \lean{win\_abs\_sum\_even}, \lean{win\_abs\_sum\_odd\_gt\_even}
in \lean{LeanPaperCWindowV1}; \lean{winValue\_odd}, \lean{winValue\_even},
\lean{winValue\_le}, \lean{winValue\_at\_last} in \lean{LeanPaperCValueV1};
\lean{count\_even\_range}, \lean{count\_odd\_range}, \lean{sum\_ite\_two},
\lean{countVal\_oddBase}, \lean{countVal\_gFamily}, \lean{countVal\_agree} in
\lean{LeanPaperCRearrangeV1}}

\paragraph{Odd $m$, lower side, and the ratios} The lower side is budget $m-1$: under
exact uniformity the last item is uniform, and the average of $|1+w|$ over $g_{m}$ is
$(m+3)/m$. S3 gives the same argument at every $\eta$. The ratios follow by arithmetic.

\leanbox{\lean{odd\_value\_zero}, \lean{odd\_uniform\_price\_ratio} in
\lean{LeanPaperCEndgameV1}}

\section{S3. The frontier (Theorems ``The even frontier'' and ``The odd frontier'')}

\paragraph{The insertion law and its marginals} Insert $(+3,-3)$ at positions $j,j+1$ of
the alternating base, $j$ uniform on $\{1,\dots,m-1\}$. Each ranking has excess exactly
$2\cdot[k=j]$, so the mixture attains $\frac{2}{D(m-1)}=V^{\mathrm{free}}_{m}$ at every budget.
With the heavy pair at position $p+1$, $p$ uniform on $\{0,\dots,m-2\}$, $+3$ lies in the
prefix of length $k$ exactly when $p<k$ and $-3$ exactly when $p<k-1$, so the class
marginals are $k/(m-1)$ and $(k-1)/(m-1)$, derived by counting rather than posited. For odd
$m$ both deviate from $k/m$ by at most $1/m$ ($\eta=1$, with the $+3$ band tight at
$k=m-1$). For even $m$ the sign flip averages them to $(2k-1)/2(m-1)$, whose deviation
$|2k-m|/2m(m-1)$ is largest at the two end budgets, where it equals $\eta_{*}/m$. Relabelling
within weight classes turns class counts into the per-item band on a law over all $m!$
permutations.

\leanbox{\lean{insRanking\_prefix}, \lean{insRanking\_excess},
\lean{insMixture\_free\_value}, \lean{insMixtureEven\_free\_value}, \lean{mixCount\_three},
\lean{mixCount\_neg\_one}, \lean{O1\_odd}, \lean{O1\_even} in \lean{LeanPaperCRankingV1};
\lean{plus\_marginal}, \lean{minus\_marginal}, \lean{plus\_deviation},
\lean{minus\_deviation}, \lean{flip\_marginal}, \lean{flip\_deviation},
\lean{flip\_deviation\_max} in \lean{LeanPaperCAttainV1};
\lean{actual\_banded\_witness\_of\_counts} in \lean{LeanPaperCRelabelBandJoinV1};
\lean{odd\_endpoint\_law}, \lean{even\_endpoint\_law} in \lean{LeanPaperCEndpointLawV1}}

\paragraph{The two counting inequalities} For every ranking of $g_{m}$:
\begin{align*}
\text{(L1), $m$ even:}\quad & r_{1}+r_{m-1}=\tfrac{2}{D}\bigl([\pi_{1}\ \text{heavy}]
  +[\pi_{m}\ \text{heavy}]\bigr),\\
\text{(L2), $m$ odd:}\quad & \textstyle\sum_{k=1}^{m-1}(|S_{k}|-b_{k})\ge
  2+2\,[\pi_{m}>0].
\end{align*}
(L1) holds because the weights sum to zero. For (L2), if $S_{k+1}=S_{k}+3$ then
$|S_{k}|+|S_{k+1}|\ge3$ while the two baselines sum to one, so the budgets straddling the
$+3$ carry excess $2$; the case analysis keeps those budgets in range, and the one case
needing its own computation is the heavy item immediately before the end.

\leanbox{\lean{two\_budget\_identity}, \lean{jump\_forces\_excess},
\lean{odd\_total\_excess}, \lean{baseline\_first}, \lean{baseline\_last} in
\lean{LeanPaperCRankingV1}; \lean{even\_total\_excess} in \lean{LeanPaperCEvenTotalV1}, \lean{hineq\_for\_breakpoint} in \lean{LeanPaperCRankingV1};
\lean{even\_rankSum} in \lean{LeanPaperCEvenRankSumV1}}

\paragraph{The band, at the two ends} The first position is the budget-$1$ prefix and the
last position is the complement of the budget-$(m-1)$ prefix, both as pointwise identities
of indicators. So $\Pr(\pi_{1}=i)\ge(1-\eta)/m$ from the band's lower side at $k=1$, and
$\Pr(\pi_{m}=i)\ge(1-\eta)/m$ from its \emph{upper} side at $k=m-1$; a one-sided band would
not give the second. On the odd side the last-position events are exclusive and summing over
the $(m-1)/2$ positive items gives $(m-1)(1-\eta)/2m$. On the even side the end events are
not exclusive and need not be: $\mathrm{big}(\pi)$ is a sum of four indicators, expectation
is additive, and the four lower bounds add to $4(1-\eta)/m$.

\leanbox{\lean{first\_eq\_inPrefix}, \lean{last\_eq\_one\_sub\_inPrefix},
\lean{endpoint\_first\_lower}, \lean{endpoint\_last\_lower}, \lean{regret\_odd\_budget},
\lean{hid\_on\_regret}, \lean{avg\_two\_le\_maxRegret}, \lean{O2\_even\_at\_ranking\_law},
\lean{O2\_even\_at\_value} in \lean{LeanPaperCEndpointV1}; \lean{hid\_at\_permutation\_law}
in \lean{LeanPaperCInstantiationV1}; \lean{first\_position\_lower},
\lean{last\_position\_lower}, \lean{band\_sum\_lower}, \lean{band\_aggregate},
\lean{exclusive\_sum\_le}, \lean{hband\_odd\_of\_prefix\_band},
\lean{hband\_even\_of\_endpoint\_band}, \lean{O2\_odd\_from\_prefix\_band},
\lean{O2\_even\_from\_endpoint\_band}, \lean{uniform\_item\_band} in \lean{LeanPaperCBandV1}}

\paragraph{Strictness below the breakpoint, with no window} With (L1), (L2) and the band
bounds, every band-feasible law satisfies, for odd $m$,
\[
\frac{\mathbb{E}_{P}[\mathrm{tot}]}{D(m-1)}\ \ge\ \frac{2}{D(m-1)}+\frac{1-\eta}{Dm},
\]
and for even $m$ the two end budgets alone exceed $V^{\mathrm{free}}_{m}$ below $\eta_{*}$.
The proofs use the free value only as an upper bound, so the margin is uniform over the
feasible set and passes to the infimum.

\leanbox{\lean{band\_floor}, \lean{end\_floor\_of\_band}, \lean{breakpoint\_lower\_even},
\lean{breakpoint\_lower\_odd}, \lean{parity\_gap\_half}, \lean{O2\_strict\_odd},
\lean{O2\_strict\_even} in \lean{LeanPaperCBreakpointV1}; \lean{odd\_breakpoint},
\lean{oddBreakpointFinish\_holds}, \lean{even\_frontier\_all} in
\lean{LeanPaperCEndpointLawV1}}

\paragraph{The even segment, and the shape} For $\eta\le\eta_{*}$ the blend of the zero-band
law (value $\frac{4}{mD}$, feasible at $0$) and the insertion law (value $V^{\mathrm{free}}_{m}$,
feasible at $\eta_{*}$) with weight $\eta/\eta_{*}$ is feasible at $\eta$ with value at most
$4(1-\eta)/mD$, since both sides are affine in $\eta$ and agree at the ends. The end budgets
give the matching lower bound. Monotonicity is nesting. The Lean proof uses explicit blends of
feasible laws rather than a general convexity lemma; convexity of $V_m$ also follows from the proved
closed form.

\leanbox{\lean{blend\_feasible}, \lean{objective\_blend\_le} in \lean{LeanPaperCEndgameV1};
\lean{value\_antitone\_of\_nested} in \lean{LeanPaperCFrontierV1}}

\paragraph{The first odd line} For odd $m$ the prefix of
length $m-1$ is $-1$ minus the last item and $m-1$ is even, so
$r_{m-1}=|1+w_{\mathrm{last}}|/D$. With $\pi_{i}=\Pr(\sigma(m-1)=i)=1-\Pr(i\in
\mathrm{prefix}_{m-1})$ the band gives $|\pi_{i}-1/m|\le\eta/m$. For any centre $a$,
$\sum\pi_{i}c_{i}=\sum(\pi_{i}-1/m)(c_{i}-a)+\frac1m\sum c_{i}$ because the deviations sum
to zero, hence $\sum\pi_{i}c_{i}\ge\frac1m\sum c_{i}-\frac\eta m\sum|c_{i}-a|$. With
$c_{i}=|1+g_{i}|\in\{2,0,2,4\}$ on the weights $-3,-1,+1,+3$ and $a=2$, the sums are $m+3$
and $m+1$. Dividing by $D$ and bounding the maximum over budgets below by budget $m-1$ gives
$V_{m}(\eta)\ge\frac{m+3-(m+1)\eta}{m(m+4)}$ for every feasible law, hence for the value.

\leanbox{\lean{split\_sum\_fn}, \lean{sum\_coef}, \lean{sum\_coef\_dev},
\lean{band\_linear\_lower}, \lean{rankOf\_isOddFamily}, \lean{regret\_last}, \lean{lastLaw},
\lean{lastLaw\_sum}, \lean{lastLaw\_band}, \lean{expect\_last}, \lean{relaxed\_value},
\lean{relaxed\_value\_odd}, \lean{relaxed\_value\_linear} in
\lean{LeanPaperCRelaxedJoinV1}}

\NOTEPAR{Status of the earlier route. \lean{relaxed\_value\_lower\_bound},
\lean{band\_expectation\_lower\_bound}, \lean{sum\_abs\_coef\_sub\_two} and
\lean{relaxed\_bound\_is\_linear\_in\_m} in \lean{LeanUnifiedStructuralRelaxedV1} prove the
same inequality for a last-position vector on a separate item type with a position band.
They were cited for Theorem~4 before \lean{LeanPaperCRelaxedJoinV1} existed, but nothing
joined them to \lean{value}. The join above replaces them as the evidence for Theorem~4.
The comparison with the free value, \lean{proved\_bound\_beats\_free} and
\lean{breakpoint\_window\_shrinks} in \lean{LeanPaperClosuresV1}, belonged to an odd-$m$
route to strictness that left a window $[(m^{2}-3)/(m^{2}-1),1)$. That route is superseded by
\lean{O2\_strict\_odd}, which leaves no window.}

\paragraph{The second odd line} Every ranking has total excess at least $2+2[\pi_{m}>0]$
by (L2). The positive items number $\frac{m-1}{2}$ ($q$ light ones and the $+3$), and the
band's upper side at $k=m-1$ puts each of them last with probability at least
$\frac{1-\eta}{m}$. So the expected total excess is at least
$2+(m-1)\frac{1-\eta}{m}$, and dividing by the $m-1$ budgets bounds the maximum below by
$V^{\mathrm{free}}_{m}+\frac{1-\eta}{mD}$. This strengthens the strictness bound above: it
counts the positive light items as well as the $+3$.

\leanbox{\lean{pos\_count}, \lean{last\_pos\_mass}, \lean{odd\_rank\_total\_pos},
\lean{odd\_line2\_law}, \lean{odd\_line2\_value} in \lean{LeanPaperCOddKinkV1}}

\paragraph{The law at the crossing point} Write $m=2s+3$ and measure masses in units
$u=\frac{1}{m^{2}(m-1)}$. The law mixes five kinds of orders, each a rearrangement of $g_{m}$
proved by matching weight counts.
\begin{itemize}
\item Version~A at slot $t$ is the insertion order $\mathrm{insWeight}(t)$.
\item Version~B at slot $t$ negates it and sets the last item to $-1$.
\item Slots $t=0,\dots,m-3$ carry $m^{2}+1$ each, half on each version, except
$m^{2}-m+4$ at $t=0$.
\item Version~A at slot $m-2$ carries $m(m-1)$.
\item The order with the light items of slot $m-2$ first and then $(-3,+3)$ carries $2$.
\item Slot $0$ with its last two light items swapped, ending $(-1,+1)$, carries $m-3$.
\end{itemize}
The masses sum to one. Every budget carries expected excess exactly $2(m^{2}+1)u$, computed
slot by slot from the single-spike profile of each version. The class counts of every prefix
lie within the band at $\eta_{t}=1-\frac{2}{m(m-1)}$, for each of the four classes. The two
light classes need $s\ge2$, so $m=3$ and $m=5$ have explicit tabulated laws, verified the
same way. A generic lemma turns ``counts, masses, band, value'' into a law feasible at
$\eta_{t}$ with that value.

\leanbox{\lean{spB}, \lean{r3W}, \lean{r1W}, \lean{cnt\_spB}, \lean{cnt\_r3W},
\lean{cnt\_r1W}, \lean{rearr\_of\_cnt}, \lean{tcount\_flip} in \lean{LeanPaperCOddKinkV1};
\lean{KIdx}, \lean{kwt}, \lean{kperm}, \lean{kq}, \lean{kq\_sum}, \lean{kex\_eq},
\lean{kslot\_excess}, \lean{kmix\_excess}, \lean{vstar}, \lean{kvalue}, \lean{pcc\_split}
in \lean{LeanPaperCOddKinkLawV1}; \lean{etaT}, \lean{slot\_heavy}, \lean{slot\_light},
\lean{band\_three}, \lean{band\_neg\_three}, \lean{band\_neg\_one}, \lean{band\_one} in
\lean{LeanPaperCOddKinkBandV1}; \lean{law\_of\_data}, \lean{kink\_law},
\lean{kink\_law\_three}, \lean{kink\_law\_five}, \lean{kink\_all} in
\lean{LeanPaperCOddFrontierV1}}

\paragraph{Assembly} The two lines meet at $\eta_{t}$, where both equal
$\frac{2(m^{2}+1)}{m^{2}(m-1)D}$. By convexity of $V_{m}$ in $\eta$, three laws suffice for
the upper side: the zero-band law, the crossing-point law and the insertion law at
$\eta=1$. Blending covers $[0,\eta_{t}]$ and $[\eta_{t},1]$, and monotonicity covers
$\eta\ge1$. The two lower lines give the other side.

\leanbox{\lean{line1}, \lean{line2}, \lean{vstar\_line1}, \lean{vstar\_line2},
\lean{etaT\_pos}, \lean{etaT\_lt\_one}, \lean{odd\_frontier\_s}, \lean{odd\_frontier},
\lean{odd\_kink\_value} in \lean{LeanPaperCOddFrontierV1}}

\NOTEPAR{Superseded evidence. Before the odd frontier was proved, the tight segment rested on
exact rational certificates at twelve points, $m\in\{5,7,9\}$ and
$\eta\in\{0,\tfrac14,\tfrac12,\tfrac34\}$ (\texttt{EXP\_paperC\_dual\_certificates}, 7/7,
\CERT), and its extent on a floating-point bisection that matched $1-\frac{2}{m(m-1)}$ to seven digits
(\NUM). Both now follow from the theorem and are kept as consistency checks only.}

\section{S4. Flow equivalence (Theorem ``Flow equivalence'')}

\paragraph{Ranking to flow} Given a law $P$ and the prefix composition at each budget, set
$x_{k}(v)=\Pr_{P}(\mathrm{comp}_{k}=v)$ and
$f_{k}(v,v')=\Pr_{P}(\mathrm{comp}_{k}=v,\mathrm{comp}_{k+1}=v')$. Conservation in both
directions is marginalisation of one joint law; the flow is non-negative and has unit mass;
and $\mathbb{E}_{P}[c(\mathrm{comp}_{k})]=\sum_{v}x_{k}(v)c(v)$ for every per-budget cost, so
the whole budget vector transports at once. The observables are named at the layer map
\lean{rankLayer}, the count vector of the first $k$ positions by type. Per-item marginals are
not functions of the composition; they become so after relabelling within each type group.

\leanbox{\lean{sum\_joint\_right}, \lean{sum\_joint\_left}, \lean{joint\_nonneg},
\lean{sum\_occ}, \lean{expectation\_via\_occ} in \lean{LeanPaperCFlowV1};
\lean{flow\_conservation\_out}, \lean{flow\_conservation\_in}, \lean{flowX\_total},
\lean{objective\_transports}, \lean{group\_count\_transports},
\lean{budget\_vector\_transports} in \lean{LeanPaperCNamingV1}; \lean{per\_item\_from\_flow}
in \lean{LeanPaperCTransport2V1}}

\paragraph{Flow to ranking: the section} A unit flow carries its own Markov law on whole
trajectories whose marginal at every budget is its occupancy; layer marginals alone would
not suffice, since a ranking is a trajectory. Every trajectory of positive mass that starts
at $0$ and ends at the multiplicity vector is the layer map of an actual permutation:
\emph{(i)} the kernel is $f/x$ where $x>0$ and $0$ otherwise, so positive mass forces every
step to raise exactly one coordinate by one; \emph{(ii)} recording which coordinate rises
gives a type sequence whose running counts equal the trajectory, by induction; \emph{(iii)} a
type sequence with the right multiplicities is a permutation, by choosing a bijection
between each pair of equal-size fibres and gluing. Uniform relabelling within each
type-preserving subgroup then makes the law invariant, supplies the per-item band from class
counts, and moves no relabelling-invariant quantity.

\leanbox{\lean{trajLaw\_marginal\_all}, \lean{relabelledLaw\_invariant},
\lean{relabelledLaw\_expectation}, \lean{linear\_objective\_from\_marginals},
\lean{kernel\_stepSupported}, \lean{trajLaw\_stepC}, \lean{matchPerm\_type} in
\lean{LeanPaperCBreakpointV1}; \lean{exists\_rank\_section}, \lean{countSec\_layer},
\lean{rankLayer\_smul} in \lean{LeanPaperCFinalV1}}

\paragraph{The domain, and the join} The feasible flows are \lean{CountFeasibleH}: a legal
one-item step at every layer below the horizon $m$ and idling after it, with unit mass. The
domain is inhabited (\lean{horizon\_flow}). The two directions are assembled as an equality
of the \emph{sets} of reachable budget vectors, so every functional of the vector has equal
infima, and the value, even frontier and odd breakpoint hold on the flow side.

\leanbox{\lean{CountFeasibleH}, \lean{horizonFlow}, \lean{horizon\_flow} in
\lean{LeanPaperCHorizonV1}; \lean{rankingFlow\_band}, \lean{bandFlowLaw\_feasible},
\lean{banded\_vectors\_agree}, \lean{actual\_banded\_vectors}, \lean{actual\_banded\_infima}
in \lean{LeanPaperCFlowBandJoinV1}; \lean{value\_eq\_flow}, \lean{even\_frontier\_flow},
\lean{odd\_breakpoint\_flow} in \lean{LeanPaperCFlowObjectiveV1}}

\NOTEPAR{Withdrawn. \lean{pathDecomposition\_count\_states} (\lean{LeanPaperCFinalV1}) is
stated over \lean{CountFeasible}, which demands a legal one-item step at \emph{every} layer.
Conservation then forces the expected number of served items to grow by one per layer while
it is bounded by $dm$, so no unit flow satisfies it (\lean{countFeasible\_empty} in
\lean{LeanPaperCVacuityV1}). That theorem therefore proves nothing and is not used. An
earlier draft of the paper cited it as the flow-to-ranking direction. The same applies to
the hypothesis-taking form \lean{pathDecomposition\_discharged}, which is not cited as
evidence. The section argument above is sound, and the join uses the repaired domain.}

\section{S5. Robustness (Theorems ``Envelope'' and ``Collapse scale'')}

\paragraph{The physical envelope} For arbitrary $g,g'$, a fixed ranking, a fixed budget and
$0<d\le\min(\|g\|_{1},\|g'\|_{1})$, the normalised regret with recomputed baseline moves by
at most $3\|g-g'\|_{1}/d$: one unit each for the prefix sum, the baseline and the
normaliser. Expectation, maximum over budgets and infimum over a fixed feasible set each
preserve the bound. The item-band feasible set does not depend on the weights, so the value
moves by the same amount. On $g_{m}$ the physical value is $V_{m}(\eta)$ (S1).

\leanbox{\lean{regret\_perturbation}, \lean{abs\_baseline\_sub\_le},
\lean{abs\_expect\_sub\_le}, \lean{abs\_sup'\_sub\_le}, \lean{abs\_inf'\_sub\_le} in
\lean{LeanPaperCFrontierV1}; \lean{normMass}, \lean{trueValue}, \lean{trueValue\_stable} in
\lean{LeanPaperCStabilityEndgameV1}; \lean{family\_value\_robust},
\lean{even\_frontier\_robust} in \lean{LeanPaperCStabilityFamilyV1}}

\paragraph{The price ratio of an actual perturbation} The two ratio clauses apply the envelope
at the two anchors $\eta=0$ (value $\frac{m+3}{m(m+4)}$, Theorem ``Zero band'') and $\eta=m$ (value
$V^{\mathrm{free}}_{m}$, since $m\ge1$ is past the odd breakpoint and the band is vacuous).
For the last clause, $\|g'\|_{1}\ge\|g_{m}\|_{1}-\varepsilon\ge m+3$ when $\varepsilon\le1$, so $d=m+3$ is
admissible and $(m+3)\cdot3\varepsilon/d=3\varepsilon$. The perturbed free value is at least
$V^{\mathrm{free}}_{m}-3\varepsilon/d>0$ because
$3/(2m(m+3))<2/((m+4)(m-1))$. Finally
$4V_{m}(0)-(m-1)V^{\mathrm{free}}_{m}=\frac{2m+12}{m(m+4)}\ge\frac{3}{2m}\ge3\varepsilon$, which is the
arithmetic of $(m-1)/4\le(V_{m}(0)-E)/(V^{\mathrm{free}}_{m}+E)$. The hypothesis is
satisfiable, since $g'=g_{m}$ meets it.

\leanbox{\lean{odd\_zero\_value}, \lean{odd\_free\_value}, \lean{normMass\_ge},
\lean{l1\_nonneg}, \lean{physical\_price\_ratio}, \lean{ratio\_arith},
\lean{physical\_price\_linear}, \lean{physical\_price\_linear\_self} in
\lean{LeanPaperCRatioJoinV1}}

\paragraph{Quantifiers in the robustness conclusion}
The universal lower bound and the one-item upper construction combine into a statement about
an infimum over a perturbation ball, not about every point in that ball. For odd $m$ define
$\mathcal R_m(\varepsilon)=\inf_{\|g'-g_m\|_1\le\varepsilon,\ V(g';m)>0}
V(g';0)/V(g';m)$. The main paper proves explicitly that, for $0<\varepsilon\le1/6$,
\begin{gather*}
\frac1{16}\min\{m-1,1/\varepsilon\}\le\mathcal R_m(\varepsilon),\\
\mathcal R_m(\varepsilon)\le\min\{(m+3)(m-1)/(2m),10/\varepsilon\}.
\end{gather*}
The wrapper is \lean{robust\_ball} in \lean{LeanPaperCRobustBallV1}. The ball price
\lean{robustPrice} is the infimum of the price ratio over the members of the ball with positive
free value. Its feasible set contains $g_m$ (\lean{family\_mem}, with \lean{family\_free\_pos}).
Every member keeps the lower bound (\lean{member\_lower}):
\lean{physical\_price\_linear} below $1/(2m)$ and \lean{collapse\_all} above it. This gives
boundedness below and \lean{robust\_lower}. The upper bound \lean{robust\_upper} applies the
infimum separately to $g_m$ (\lean{family\_price}) and to the one-item construction at
$\delta=\varepsilon$ (\lean{witness\_mem}). No minimiser of the ball problem is required.
Positive scaling, $g'=(1+\varepsilon/(m+4))g_m$, preserves all normalised regrets
(\lean{trueValue\_smul}). It is a member at distance exactly $\varepsilon$ that keeps the
unperturbed price (\lean{rescaled\_keeps\_price}), a counterexample to the stronger, incorrect
statement that every perturbation beyond $1/m$ forces the ratio to decay.

\paragraph{The collapse, upper side} Move light item $1$ from $-1$ to $-1+\delta$,
$0<\delta\le1$. Then $\|g_{m}-g'\|_{1}=\delta$ and $\|g'\|_{1}=m+4-\delta$. At odd budgets
the best drop takes the moved item and costs at most $1-\delta$. At even budgets the best
drop leaves it out and costs $0$. An order pays at least $\delta$ at an odd budget whose
prefix excludes the moved item, because every other prefix sum there is an odd integer. It
also pays at least $\delta$ at an even budget whose prefix includes it, because that sum is an
even integer plus $\delta$. Each pair $(2i+1,2i+2)$ contains one such budget, so the total
over budgets is at least $\delta\frac{m-1}{2}$. Hence the free value is at least
$\frac{\delta}{2\|g'\|_{1}}>0$. The envelope bounds the zero-band value by
$V_{m}(0)+\frac{3\delta}{d}$, and the arithmetic gives a ratio of at most
$\frac{10}{\delta}$.

\leanbox{\lean{pI}, \lean{gD}, \lean{l1\_gD}, \lean{normMass\_gD}, \lean{baseline\_odd},
\lean{baseline\_even}, \lean{prefix\_gD}, \lean{pays}, \lean{raw\_ge\_pays},
\lean{pays\_pairs}, \lean{objective\_gD\_ge}, \lean{value\_gD\_ge}, \lean{collapse\_upper},
\lean{collapse\_arith} in \lean{LeanPaperCCollapseScaleV1}}

\paragraph{The collapse, lower side} Let $\varepsilon=\|g_{m}-g'\|_{1}\in[\frac1m,\frac16]$
with a positive unconstrained value. The envelope with $d=m+4-\varepsilon$ bounds the
zero-band value below by $V_{m}(0)-\frac{3\varepsilon}{d}$ and the free value above by
$V^{\mathrm{free}}_{m}+\frac{3\varepsilon}{d}$. On this window both are dominated by the
$\varepsilon$ terms up to constants, and the ratio is at least $\frac{1}{16\varepsilon}$. The
positivity hypothesis is displayed because the envelope's lower bound on the free value can be
negative here.

\leanbox{\lean{collapse\_lower} in \lean{LeanPaperCCollapseScaleV1}}

\paragraph{The window $(\frac{1}{2m},\frac1m)$} The same envelope with $d=m+3$ applies to every
$g'$ with $\varepsilon\le\frac1m$, since then $\frac{3\varepsilon}{m+3}\le\frac{3}{m(m+3)}$ stays
below the zero-band value. With $x=m$ and $y=\frac{3\varepsilon}{x+3}$, the ratio bound
$\frac{a-y}{f+y}$ is at least $\frac{x-1}{8}$. That reduces to
$x(x+4)(3x^{2}+9x-12)\ge0$. With the lower side above, the ratio is at least
$\frac1{16}\min(m-1,\frac1\varepsilon)$ on the whole range $0<\varepsilon\le\frac16$. The
unperturbed family meets the hypotheses.

\leanbox{\lean{window\_arith}, \lean{collapse\_window}, \lean{collapse\_all},
\lean{collapse\_window\_satisfiable} in \lean{LeanPaperCCollapseWindowV1}}

\NOTEPAR{Status of the fixed-normaliser floor. \lean{price\_ratio\_robust}
(\lean{LeanPaperCBreakpointV1}), \lean{ratio\_lower\_of\_perturbation},
\lean{ratio\_linear\_at\_small\_eps}, \lean{ratio\_bounded\_at\_fixed\_eps}
(\lean{LeanPaperCTransport2V1}) and \lean{ratio\_collapse\_lower}
(\lean{LeanPaperCCollapseV1}) are arithmetic. They take the perturbed zero-band and free
values as hypotheses, with the normaliser held at $D=m+4$, and pin the resulting floor
expression between $\frac18\min(m-1,1/\varepsilon)$ and $2/\varepsilon$. No declaration
instantiates their value hypotheses at an actual perturbed program, so they bound a floor
expression and not the price. They are drawn as the floor lines of the stability figure and
are not cited for any theorem of the main text. \lean{LeanPaperCRatioJoinV1} is the
instantiated version, for the physical program with the moving normaliser.}

\paragraph{The search} Random restarts followed by coordinate descent on the price ratio,
over the $m!$ permutation formulation, with baseline and normaliser recomputed at every
step, at $m=5,6,7$ and $\varepsilon\le1.6$, over three families: adversarial, shrinking the
heavy anchors, and a uniform shift (\texttt{EXP\_paperC\_robustness}, 4/4). A minimisation
returns attained ratios, upper bounds on the searched infimum, so it can exhibit a collapse but
not exclude one. (\NUM)

\section{S6. The ex-post values (Proposition ``Free ex-post value'' and Theorem ``Banded
ex-post value'')}

Every ranking of $g_{m}$ has an interior budget with excess at least $2$: one heavy item is
not last, a zero-excess budget has $|S_{k}|\le1$, a step of $3$ cannot join two such budgets,
and excess is even. The flat ranking (even $m$) and the base ranking (odd $m$) have excess
at most $2$ everywhere. So the free ex-post value is $2/D$, and $\frac{2}{D}\big/\frac{2}{D(m-1)}=m-1$.

\leanbox{\lean{exists\_defect}, \lean{worstBudget\_even}, \lean{worstBudget\_odd},
\lean{pathValue\_eq}, \lean{path\_over\_free}, \lean{objective\_le\_pathObjective} in
\lean{LeanPaperCExPostV1}}

\paragraph{The banded ex-post value} Let $V^{\mathrm{post}}_{m}(\eta)$ be the infimum of
$\mathbb{E}_{P}[\max_{k}r_{k}]$ over band-feasible laws.
\begin{itemize}
\item \emph{Lower side.} Every order has worst excess at least $2$. An odd order with $+3$
last pays $4$ at budget $m-1$, since the prefix sum there is $-4$ and the baseline is $0$. The
band puts $+3$ last with probability at least $\frac{1-\eta}{m}$.
\item \emph{Upper side.} A uniform law on the $m$ cyclic shifts of one order has marginals
exactly $\frac{k}{m}$, so it is feasible at $\eta=0$. The window sum of a rotation is read off
the ordinary prefix sequence. Every rotation of the flat even order has worst excess at most
$2$. Every rotation of the odd base order has worst excess at most $2$, except the one that
ends with $+3$, which has at most $4$. The insertion mixture has worst excess $2$ on every
order and is feasible at $\eta=1$. Blending the rotation law with it covers $[0,1]$.
\item \emph{Symmetrisation.} Relabelling within weight classes preserves the worst budget
of every order, so class-count feasibility gives per-item feasibility.
\end{itemize}

\leanbox{\lean{preG}, \lean{winG}, \lean{winG\_eq}, \lean{rot\_prefix}, \lean{preN\_small},
\lean{odd\_rot\_excess}, \lean{flatPre\_mid}, \lean{flatPre\_total}, \lean{flatPre\_small},
\lean{even\_rot\_excess}, \lean{postValue}, \lean{post\_le}, \lean{post\_ge},
\lean{pathObjective\_blend}, \lean{post\_ge\_two}, \lean{worstBudget\_plus},
\lean{post\_ge\_odd}, \lean{rankOf\_wperm}, \lean{worstBudget\_wperm}, \lean{sym\_sum},
\lean{insertion\_post}, \lean{odd\_rotation\_post}, \lean{even\_rotation\_post},
\lean{expost\_even}, \lean{expost\_odd} in \lean{LeanPaperCExPostBandV1}}

\section{S7. Any number of types (Theorems ``The flow program for every typed multiset'' and
``The free threshold'', Proposition ``The additive price is stable'')}

\paragraph{Typed flows} The core of the four-type flow equivalence, the agreement of banded
vectors between rankings and count-feasible flows, holds for any type map and any cost on
count states. It is instantiated at the physical cost of a typed multiset. On a ranking the prefix
sum is the weighted composition count, so the count cost is the normalised regret exactly.
\leanbox{\lean{typedG}, \lean{typedCost}, \lean{prefixSet\_sum\_eq}, \lean{typedCost\_eq},
\lean{typed\_flow\_equivalence}, \lean{typed\_banded\_infima}, \lean{card\_state},
\lean{family\_is\_typed} in \lean{LeanPaperCTypedFlowV1}}

\paragraph{The box} The count-feasibility condition bounds every occupied state by the
multiplicities. The box is a product of initial segments of $\{0,\dots,m\}$.
\leanbox{\lean{mult\_le}, \lean{box}, \lean{support\_box}, \lean{card\_box} in
\lean{LeanPaperCFreeBandV1}}

\paragraph{The free threshold} The band-feasible set is closed and inside the unit ball, hence
compact, and the objective is a finite maximum of linear forms. So the value is attained, and
so is the minimum deviation over the compact set of free optima.
\begin{itemize}
\item If $V_g(\eta)=V^{\rm free}_g$, an optimal law at $\eta$ is a free optimum of deviation at
most $\eta$.
\item Conversely, a free optimum of deviation $\eta_{\rm free}\le\eta$ is feasible at $\eta$.
\end{itemize}
The flow form rewrites $V_g$ as the infimum of the same functional over the flow side.
\leanbox{\lean{marg}, \lean{dev}, \lean{feasible\_iff\_dev}, \lean{feasible\_m\_eq},
\lean{isCompact\_feasible}, \lean{attain\_value}, \lean{Vfree}, \lean{OptFree}, \lean{etaFree},
\lean{attain\_eta}, \lean{free\_threshold}, \lean{flowValue}, \lean{trueValue\_eq\_flow},
\lean{free\_threshold\_flow} in \lean{LeanPaperCFreeBandV1}}

\paragraph{The additive price} Stability of $\Delta_g$ is two applications of the value
Lipschitz bound, at $\eta$ and at $\eta=m$.
\leanbox{\lean{gap}, \lean{gap\_nonneg}, \lean{gap\_antitone}, \lean{gap\_zero\_iff},
\lean{gap\_stable} in \lean{LeanPaperCFreeBandV1}}

\paragraph{One item of each type} The baselines are $b,0,b$ at budgets $1,2,3$: a single item costs at
least $b$, the $a$-pair cancels, and a triple is minus its missing item. The prefixes of an order
are $\{\sigma_0\}$, $\{\sigma_0,\sigma_1\}$ and the complement of $\{\sigma_3\}$. The six
orders are given by their tables and inverses. Band feasibility is checked at every $k\le4$ and
every item, and the budget costs are evaluated order by order.
\leanbox{\lean{g4}, \lean{abs\_g4}, \lean{normMass\_g4}, \lean{sum\_compl\_single},
\lean{prefix\_one}, \lean{prefix\_two}, \lean{prefix\_three}, \lean{raw\_one}, \lean{raw\_two},
\lean{raw\_three}, \lean{marg\_one}, \lean{mkPerm}, \lean{law\_sum}, \lean{raw\_values} in
\lean{LeanPaperCTwoMagFourV1}}

\paragraph{The two-magnitude map} \texttt{scripts/EXP\_paperC\_typed\_frontier.py} (source
SHA-256 prefix \texttt{2a572a64ce84}) wrote
\texttt{research/unified\_20260919/paperC\_typed\_frontier.json}. It has $240$ count vectors,
$4$ ratios and $41$ band values, and $3/3$ self-checks. HiGHS floating point. The refinement
\texttt{scripts/EXP\_paperC\_typed\_transition.py} reuses that solver and wrote
\texttt{paperC\_typed\_transition.json}. It covers counts $(1,1,2,2)$ and $(2,2,2,2)$, $21$
ratios in $[2,3]$ and $49$ band values, with $3/3$ self-checks, one of which reproduces the
stored rows at $a/b=2,3$. \NUM{}

\section{S8. Knowing the budget, and recourse (Proposition ``Knowing, editing, committing'',
Theorems ``Knowing the budget is worth nothing on a band range'' and ``Four items'')}

\paragraph{Objects} A subset law at budget $k$ is a probability vector on subsets, supported on
$k$-sets, whose item marginals lie within $\eta/m$ of $k/m$. Its cost is
$\sum_SQ(S)(|\sum_{i\in S}g_i|-b_k)/\|g\|_1$. The layer value is the infimum of the cost, and the
budget-aware value is its maximum over the budgets. A recourse plan at budget $k$ is a
nonnegative $J(\sigma,T)$ with $\sum_TJ(\sigma,T)=P(\sigma)$, supported on pairs with $|T|=k$ and
$|S_k(\sigma)\setminus T|\le R$, whose final item marginals lie in the band. The recourse value
is the infimum, over band-feasible $P$ and plans at each budget, of the worst budget's final cost.
\leanbox{\lean{ind}, \lean{SubFeasible}, \lean{subCost}, \lean{layerValue}, \lean{awareValue},
\lean{PlanFeasible}, \lean{planCost}, \lean{recourseSet}, \lean{recourseValue} in
\lean{LeanPaperCBudgetAwareV1}}

\paragraph{The sandwich} The push of $P$ to its budget-$k$ prefix, $S\mapsto\Pr_P(S_k=S)$, has the
same marginals and the same cost; so every layer value is at most $\mathbb E_Pr_k$, and
$B_g\le V_g$. The final-set law of a plan is a subset law in the band, so $B_g\le V^{(R)}_g$. The
plan that keeps the prefix is feasible for every $R$, so $V^{(R)}_g\le V_g$. At $R=0$ a
supported $T$ contains the prefix and has its size, so it is the prefix, and every plan has the
cost of its base law: $V^{(0)}_g=V_g$. A plan for $R$ is a plan for $R'\ge R$. For the last
statement, two $k$-sets $S,T$ have $|S\setminus T|\le|S|=k$ and $S\setminus T\subseteq T^{c}$,
so $|S\setminus T|\le m-k$. Draw the final set independently from a subset law within
$\varepsilon$ of each layer's infimum; the plan is feasible when $\min(k,m-k)\le R$, and its
cost is that law's cost. Let $\varepsilon\to0$.
\leanbox{\lean{subCost\_nonneg}, \lean{cost\_nonneg}, \lean{layer\_le}, \lean{layer\_nonneg},
\lean{pushK}, \lean{push\_sum}, \lean{inPrefix\_eq\_ind}, \lean{push\_feasible}, \lean{push\_cost},
\lean{layer\_le\_obj}, \lean{sub\_nonempty}, \lean{aware\_le\_true}, \lean{finalLaw},
\lean{final\_feasible}, \lean{layer\_le\_plan}, \lean{idPlan}, \lean{idPlan\_feasible},
\lean{idPlan\_cost}, \lean{aware\_le\_recourse}, \lean{recourse\_le\_true}, \lean{recourse\_zero},
\lean{recourse\_mono}, \lean{sdiff\_card\_le}, \lean{indepPlan}, \lean{indep\_feasible},
\lean{indep\_cost}, \lean{recourse\_full} in \lean{LeanPaperCBudgetAwareV1}}

\paragraph{Linear certificates} If $\alpha+\sum_{i\in S}\beta_i\le f(S)$ on $k$-sets, then for a
subset law $Q$ in the band,
$\sum_SQ(S)f(S)\ge\alpha+\sum_i\beta_i\Pr_Q(i\in S)$, and each term obeys
$\beta_i\Pr_Q(i\in S)\ge\beta_ik/m-|\beta_i|\eta/m$.
\leanbox{\lean{linear\_lower}, \lean{layer\_ge\_of\_linear} in \lean{LeanPaperCBudgetAwareV1}}

\paragraph{The family} The upper side is $B\le V$ and the frontiers of Theorems ``The even
frontier'' and ``The odd frontier''. On the stated ranges the frontier is its first line: in even
size $\frac{2}{(m+4)(m-1)}\le\frac{4(1-\eta)}{m(m+4)}$ exactly when $\eta\le\eta_*$, and in odd size
the second line lies below the first exactly when $\eta\le\eta_t$. The lower side is one layer.
\begin{itemize}
\item Even $m$, budget one: $\beta_i=|g_i|-1$ and $\alpha=0$, an equality on singletons since
$b_1=1$. The family sums give $\sum\beta_i=\sum|\beta_i|=4$, so the layer is at least
$4(1-\eta)/(m(m+4))$.
\item Odd $m$, budget $m-1$: a $k$-set is $\{j\}^{c}$, its sum is $-1-g_j$ and $b_{m-1}=0$. With
$\beta_i=2-|-1-g_i|$ and $\alpha=5-m$, the certificate is an equality on every such set. The family
sums give $\sum\beta_i=m-3$ and $\sum|\beta_i|=m+1$, so the layer is at least
$(m+3-(m+1)\eta)/(m(m+4))$.
\end{itemize}
\leanbox{\lean{aware\_even}, \lean{aware\_odd} in \lean{LeanPaperCBudgetGmV1};
\lean{sum\_fn\_gFamily} in \lean{LeanPaperCRelaxedJoinV1};
\lean{familyBaselineIdentification\_holds} in \lean{LeanPaperCBaselineAttainV1}}

\paragraph{Four items} Point-mass laws on explicit sets give the upper sides: singletons with
weights $(h,h,\frac12-h,\frac12-h)$, $h=\max(1-\eta,0)/4$, at budget one; their complements at
budget three; the two cancelling pairs at budget two. The lower side at budget one is the
certificate $\beta=(a-b,a-b,0,0)$. The plan uses the base law of the four cyclic orders, pushed
to permutations, which is feasible at band $0$. At budgets one and three the final set is drawn
independently. At budget two a kernel keeps a cancelling prefix and sends a mixed one to each
cancelling pair with probability $\frac12$. The cyclic prefixes of size two are $\{0,1\}$,
$\{1,2\}$, $\{2,3\}$, $\{3,0\}$, so each final pair has probability $\frac12$ and every item
marginal is $\frac12$. A mixed pair of size two differs from each cancelling pair in one item, by
exhaustion over the subsets of four items.
\leanbox{\lean{ptLaw}, \lean{ptLaw\_sum}, \lean{ptLaw\_feasible}, \lean{kerPlan},
\lean{kerPlan\_feasible}, \lean{kerPlan\_cost}, \lean{Q1\_feasible}, \lean{Q2\_feasible},
\lean{Q3\_feasible}, \lean{Q1\_cost}, \lean{Q2\_cost}, \lean{Q3\_cost}, \lean{lay1\_lower},
\lean{aware\_twoMag}, \lean{advantage\_twoMag}, \lean{P0\_feasible}, \lean{P0\_expect},
\lean{K2\_sum}, \lean{mixed\_close}, \lean{K2\_supp}, \lean{K2\_band}, \lean{K2\_cost},
\lean{recourse\_one\_twoMag}, \lean{recourse\_zero\_twoMag}, \lean{exchange\_needed} in
\lean{LeanPaperCBudgetFamiliesV1}}

\paragraph{The numerical panel} \NUM{} \texttt{scripts/EXP\_paperC\_budget\_recourse.py} solves each
layer as a linear program over count vectors with the type-marginal band, and the recourse value
as the count flow of Theorem ``The flow program for every typed multiset'' plus, at each interior
budget, a transport between count vectors of radius $R$ with the band and the regret on its
output. It refuses to write unless the four-item formulas, the two family ranges, the chain
$B\le V^{(R+1)}\le V^{(R)}\le V$ and $V^{(\lfloor m/2\rfloor)}=B$ all hold to $10^{-7}$. The run
(\texttt{paperC\_budget\_recourse.json}, $136$ s) covers $g_m$ for $m=4,\dots,10$, all
$R\le\lfloor m/2\rfloor$, and counts $(1,1,2,2)$ and $(2,2,2,2)$ at $a/b\in\{1.5,2,2.6,3,5\}$.
\begin{itemize}
\item On $g_m$ the first grid point with $A>10^{-6}$ is the first one past each proved range:
$0.34,0.91,0.41,0.96,0.43,0.98,0.45$ for $m=4,\dots,10$.
\item The least $R$ with $V^{(R)}=B$ is $0$ up to that point and $1$ after it, on every
instance.
\item With repeated types, the first gap is at $\eta=0$ for $a/b\le2$, near $0.26$--$0.29$ at
$a/b=2.6$, and near $0.41$--$0.43$ at $a/b\ge3$.
\end{itemize}
Symmetrising a subset law within types does not change its cost and keeps its band, so the
count layer has the item layer's value, and a count transport lifts to item exchanges by the same
relabelling applied to the ranking and the final set; both reductions are proved
(Proposition ``The count model is exact'' of the main text), so the recourse column is the item
model's value up to the solver's floating point.
\leanbox{\lean{symSub\_feasible}, \lean{symSub\_cost}, \lean{push\_countFeasible},
\lean{lift\_feasible}, \lean{layer\_costs\_eq\_count} in \lean{LeanPaperCCountReductionV1};
\lean{exists\_nested}, \lean{lift\_plan}, \lean{proj\_transport} in \lean{LeanPaperCCountRecourseV1}}

\section{S9. Build record}

The corpus is verified by \texttt{verify\_all.sh}, which freezes its inputs, builds every
listed module and every local import from an empty cache in import order, counts
\texttt{sorry} and \texttt{native\_decide} outside comments, checks every audited axiom
closure against the whitelist, and compares the frozen bytes with the working directory at
the end. The gate records per-module source and olean hashes. The last full gate that passed,
\texttt{full6\_20260925}, built $214$ modules and audited $1666$ theorems, with no
\texttt{sorry}, no \texttt{native\_decide}, the standard axioms, and frozen sources identical to
the working copy afterwards. It includes every module this supplement cites except two.
\lean{LeanUnifiedStructuralRelaxedV1} imports only Mathlib and passes the smoke gate
\texttt{smokeUNI\_20260925} ($2$ modules, $16$ theorems). \lean{LeanPaperCCollapseWindowV1} was
written after it. It compiles with the standard axioms and no \texttt{sorry}, and it is in the
module list of the full gate \texttt{full7\_20260925}. Before that full gate, the modules
added after \texttt{full5\_20260924} ($197$ modules, $1553$ theorems) were covered by smoke gates
over their full import closures:
\begin{itemize}
\item \texttt{smokeCJ\_20260924}: \lean{LeanPaperCRelaxedJoinV1} and
\lean{LeanPaperCRatioJoinV1}; $49$ modules, $449$ theorems.
\item \texttt{smokeCOF\_20260924}: \lean{LeanPaperCOddFrontierV1}; $51$ modules, $459$
theorems.
\item \texttt{smokeCC3\_20260924}: \lean{LeanPaperCCollapseScaleV1} and
\lean{LeanPaperCExPostBandV1}; $52$ modules, $463$ theorems.
\end{itemize}
The modules of Section~S7 were written after \texttt{full7\_20260925}. Each passes a smoke gate
over its full import closure, with the standard axioms, no \texttt{sorry}, no
\texttt{native\_decide}, and no drift:
\begin{itemize}
\item \texttt{smokeUP1\_20260925}: \lean{LeanPaperCTypedFlowV1}; $79$ modules, $762$ theorems.
\item \texttt{smokeUP2\_20260925}: \lean{LeanPaperCFreeBandV1}; $107$ modules, $934$ theorems.
\item \texttt{smokeUP3\_20260925}: \lean{LeanPaperCTwoMagFourV1}; $106$ modules, $924$ theorems.
\item \texttt{smokeUP4\_20260925}: \lean{LeanPaperCBudgetAwareV1}, \lean{LeanPaperCBudgetFamiliesV1},
\lean{LeanPaperCBudgetGmV1}; $139$ modules, $1199$ theorems.
\item \texttt{smokeC1\_20260926}: \lean{LeanPaperCTwoMag2211V1}, \lean{LeanPaperCBudgetSymV1},
\lean{LeanPaperCBudgetSharpV1}, \lean{LeanPaperCCountReductionV1}, \lean{LeanPaperCCountRecourseV1};
$65$ modules, $607$ theorems.
\end{itemize}

\end{document}
